\subsection{\texorpdfstring{Bounded complex measures \(^{6}\)}{Bounded complex measures (note 6)}}

For every complex measure \(m\) on \(\mathrm{T}\) , one sets

\[
\| m \| = \sup _ {\| f \| \leqslant 1,   f \in \mathcal {K} _ {\mathbf {C}} (\mathrm{T})} | m (f) |.
\]

One says that m is bounded if \(\|m\| < +\infty\) ; it comes to the same to say that m is continuous on \(\mathcal{K}_{\mathbf{C}}(\mathrm{T})\) equipped with the topology of uniform convergence, hence can be extended to a continuous linear form (of norm \(\|m\|\) ) on the Banach space \(\overline{\mathcal{K}_{\mathbf{C}}(\mathrm{T})}\) of continuous complex functions tending to 0 at infinity.

Lemma 5. --- Let m be a complex measure on T, f an m-integrable complex function. Then \(\int |f| d|m| = \sup \left| \int f h dm \right|\) , as h runs over the set of functions in \(\mathcal{K}_{\mathbf{C}}(\mathbf{T})\) such that \(|h(t)| \leqslant 1\) for all \(t \in T\) .

If \(m = g \cdot |m|\) , then \(\int |f| d|m| = \int |fg| d|m|\) and \(\int f h dm = \int f gh d|m|\) . Set \(\zeta(t) = 0\) when \(f(t)g(t) = 0\) , and \(\zeta(t) = \frac{f(t)g(t)}{|f(t)g(t)|}\) when \(f(t)g(t) \neq 0\) ; \(\zeta\) is \(|m|\) -measurable, thus for every \(\varepsilon > 0\) there exists a compact subset K of T such that \(\int_{T-K} |f| d|m| \leqslant \varepsilon\) , the restriction of \(\zeta\) to K is continuous, and \(|\zeta(t)| = 1\) on K. Therefore, by virtue of Urysohn's theorem, there exists a continuous function \(\zeta_1\) defined on T, with complex values, such that \(\zeta_1 = \zeta\) on K and such that \(|\zeta_1(t)| \leqslant 2\) and \(\zeta_1(t) \neq 0\) for every \(t \in T\) ; setting \(h(t) = \zeta_1(t)/|\zeta_1(t)|\) , we see that h is continuous on T, coincides with \(\zeta\) on K, and is such that \(|h(t)| = 1\) for all \(t \in T\) . Finally, let u be a continuous mapping of T into {[}0,1{]}, equal to 1 on K and with compact support; setting \(h_1 = h^{-1}u\) , we have

\[
\left| \int f h _ {1} d m - \int | f | d | m | \right| \leqslant 2 \int_ {\mathrm{T-K}} | f | d | m | \leqslant 2 \varepsilon ,
\]

which proves the lemma.

PROPOSITION 13. --- Let m be a complex measure, and \(\mu = |m|\) . For m to be bounded, it is necessary and sufficient that \(\mu\) be bounded, in which case \(\|m\| = \|\mu\|\) .

We have \(m = g \cdot \mu\) , where \(g\) is \(\mu\) -measurable and \(|g(t)| = 1\) for all \(t \in \mathbf{T}\) . If \(\mu\) is bounded then, for every function \(f \in \mathcal{H}_{\mathbf{C}}(\mathbf{T})\) ,

\[
| m (f) | = \left| \int f g d \mu \right| \leqslant \mathrm{N} _ {\infty} (f g) \| \mu \| = \| f \| \cdot \| \mu \|,
\]

therefore m is bounded and \(\|m\|\leqslant\|\mu\|\) . If m is bounded, then for every \(f\in\mathcal{K}_{\mathbf{C}}(\mathrm{T})\) we have, on taking into account Lemma 5,

\[
| \mu (f) | \leqslant \| f \| \cdot \| m \|,
\]

therefore \(\mu\) is bounded and \(\| \mu \| \leqslant \| m\|\) . Whence the proposition.

COROLLARY. --- Let m be a bounded complex measure. Every function \(\mathbf{f} \in \mathcal{L}_{\mathrm{F}}^{\infty}(m)\) is then m-integrable, and \(\left|\int f dm\right| \leqslant N_{\infty}(\mathbf{f}) \|m\|\) .

For, \(\mathbf{f}\) is \(m\) -measurable and, setting \(\mu = |m|\) , we have

\[
\int^ {*} | \mathbf {f} | d \mu \leqslant \mathrm{N} _ {\infty} (\mathbf {f}) \| \mu \| = \mathrm{N} _ {\infty} (\mathbf {f}) \| m \|,
\]

therefore \(\mathbf{f}\) is \(|m|\) -integrable (Ch. IV, §5, No.~6, Th. 5) and

\[
\left| \int \mathbf {f} d m \right| \leqslant \int | \mathbf {f} | d \mu \leqslant \mathrm{N} _ {\infty} (\mathbf {f}) \| m \|.
\]

\begin{enumerate}
\def\labelenumi{\arabic{enumi}.}
\setcounter{enumi}{9}
\tightlist
\item
  Image of a complex measure; induced complex measure; product of complex measures \(^{7}\)
\end{enumerate}

Let \(m\) be a complex measure on \(T\) , and let \(\pi\) be a mapping of \(T\) into a locally compact space \(X\) . We shall say that \(\pi\) is \(m\) -proper if \(\pi\) is \(|m|\) -proper (Ch. V, §6, No.~1, Def. 1); it is then immediate that for every function \(f \in \mathcal{H}_{\mathbf{C}}(X)\) , the function \(f \circ \pi\) is essentially \(m\) -integrable and

\[
\left| \int (f \circ \pi) d m \right| \leqslant \int | f \circ \pi | d | m | = \int | f | d (\pi (| m |)),\tag{8}
\]

therefore the mapping \(f \mapsto \int (f \circ \pi) dm\) is continuous on \(\mathcal{K}_{\mathbf{C}}(\mathbf{X})\) , in other words is a complex measure on \(\mathbf{X}\) , which is denoted \(\pi(m)\) and called the image of \(m\) under \(\pi\) . Moreover, it follows from (8) that \(|\pi(m)| \leqslant \pi(|m|)\) . If \(m\) and \(m'\) are two complex measures on \(T\) and if \(\pi\) is both \(m\) -proper and \(m'\) -proper, then \(\pi\) is \((\lambda m + \lambda'm')\) -proper for any complex scalars \(\lambda, \lambda'\) , and \(\pi(\lambda m + \lambda'm') = \lambda\pi(m) + \lambda'\pi(m')\) .

Let Y be a locally compact subspace of T. For every function \(f \in \mathcal{K}_{\mathbf{C}}(\mathbf{Y})\) , the function \(f'\) on T, defined by \(f'(t) = f(t)\) if \(t \in Y\) and by \(f'(t) = 0\) if \(t \notin Y\) , is m-integrable (Ch. IV, §5, No.~7); it is immediate that the mapping \(f \mapsto \int f' dm\) is a complex measure on Y, called the measure induced on Y by m and denoted \(m_{Y}\) . If \(m = g \cdot |m|\) , it is clear that \(m_{Y} = g_{Y} \cdot |m|_{Y}\) , where \(g_{Y}\) is the restriction to Y of the function g, which is locally integrable for \(|m|_{Y}\) (Ch. V, §7, No.~1); moreover, since \(|g_{Y}| = 1\) locally almost everywhere for \(|m|_{Y}\) (Ch. V, §7, No.~1, Cor. 1 of Prop. 1), we have \(|m_{Y}| = |m|_{Y}\) .

Let T and T' be two locally compact spaces, m (resp. \(m'\) ) a complex measure on T (resp. T'). Write \(m = g \cdot |m|\) and \(m' = g' \cdot |m'|\) . The function \(g \otimes g'\) is locally integrable on T × T' for the positive measure \(|m| \otimes |m'|\) (Ch. V, §8, No.~5, Prop. 10), and one verifies immediately that if g (resp. \(g'\) ) is replaced by a function \(g_1\) (resp. \(g'_1\) ) equal to g (resp. \(g'\) ) locally almost everywhere for \(|m|\) (resp. \(|m'|\) ), then \(g_1 \otimes g'_1\) is equal to \(g \otimes g'\) locally almost everywhere for \(|m| \otimes |m'|\) . The complex measure \((g \otimes g') \cdot (|m| \otimes |m'|)\) on T × T' therefore depends only on m and \(m'\) ; it is denoted \(m \otimes m'\) and is called the product measure of m and \(m'\) . Since \(|g \otimes g'| = 1\) locally almost everywhere for \(|m| \otimes |m'|\) (Ch. V, §8, No.~2, Prop. 4), we have \(|m \otimes m'| = |m| \otimes |m'|\) .

The reader will easily verify that all of the propositions proved in Ch. V relative to the image of a positive measure, the measure induced by a positive measure, and the product of positive measures, except those in which upper integrals or essential upper integrals intervene, remain valid when the positive measures are replaced by arbitrary complex measures.

Finally, one defines as in §1 the concept of scalarly essentially m-integrable function for a complex measure m; for a function f to have this property, it is necessary and sufficient that f be scalarly essentially integrable with respect to \(|\mu_{1}|\) and \(|\mu_{2}|\) , where \(\mu_{1}\) and \(\mu_{2}\) are the real and imaginary parts of m, in which case \(\int f dm = \int f d\mu_{1} + i \int f d\mu_{2}\) . We leave to the reader the task of carrying over the results of §1 to complex measures.

